\begin{lemma}[The actual pulled-canonical trace under its fixed native chart comparisons]
\label{lem:kn-regular-pulled-canonical-fixed-chart-restriction}
\lean{EnvelopeProperty.Vanishing.NadelGlobalCanonicalTrace.regularPulledCanonicalTrace_wholeChart_fixed_restriction_captured}
\uses{lem:kn-original-line-trace-native-fixed-restriction,lem:kn-original-whole-chart-structure-square,def:kn-canonical-whole-chart-fixed-native-comparisons,def:nadel-regular-pulled-canonical-trace}
\leanok
The SAME actual global regular pulled-base canonical trace restricts to the
actual whole-chart line trace under the FIXED native base and pulled-line
isomorphisms. The latter trace is the original inverse projection formula
applied to the original whole coefficient-algebra trace on the chart's native
canonical base line. No chart agreement or choice of compatible isomorphisms
is a premise. The formal ordinary equality records these same actual maps
after unfolding their stored wrappers.
\end{lemma}
\begin{proof}
Apply the original fixed-line restriction theorem to the actual original
cover and whole chart. The already proved global structure-trace chart square
supplies its entire scalar equation internally. Specify the base line
isomorphism to be the actual open canonical differential; it is invertible
because that base map is an open immersion. Pullback preserves the original
canonical line's invertibility, and this same differential transports it to
the chart's canonical line. The resulting source and target comparisons are
literally the two fixed native chart isomorphisms. Unfold only these stored
comparisons and the original regular pulled-line trace.
\end{proof}
